\SourceLeaf{357}

\noindent wenn substituirt wird $X = m'm''$, $X' = mm''$, $X'' = mm'$, während der sechsfache
Cubikinhalt der Pyramide, die sich dadurch mit dem Nullpunkte bildet, $= mm'm''\BareRoot{D}$
wird; folglich [ist] das Perpendikel
\[
= \BareRoot{\dfrac{D}{F\left(\dfrac{1}{m}, \dfrac{1}{m'}, \dfrac{1}{m''}\right)}};
\]

$T$, $T'$, $T''$ beziehen sich eben so auf die Form $\begin{pmatrix} AD, A'D, A''D \\ BD, B'D, B''D \end{pmatrix}$
wie $t$, $t'$, $t''$ auf $\begin{pmatrix} a, a', a'' \\ b, b', b'' \end{pmatrix}$.



\begin{center}
\bfseries BEMERKUNG.
\end{center}

Die vorstehende Notiz befindet sich in einem Handbuche und stammt aus dem Juli 1831. Da sie
zeigt, dass \textsc{Gauss} sich auch mit complexen Grössen von mehr als zwei Einheiten beschäftigt
hat, die er bereits in der Schlussbemerkung seiner Selbstanzeige der Theoria residuorum
biquadraticorum, Commentatio secunda vom 23. April 1831 (Bd. II. S. 169 bis 198) erwähnt
hatte, so hat sie hier noch einmal Platz gefunden, obwohl sie schon in Bd. II. S. 305 bis
307 wegen ihrer Beziehung zur Theorie der ternären Formen abgedruckt ist.
\begin{flushright}
\textsc{Stäckel.}
\end{flushright}


